\begingroup
\small
\setlength{\parskip}{0pt}
\begin{itemize}
  \setlength{\itemsep}{0.25em}
  \raggedright
  \item Cátedra de Arquitectura de Computadoras, \textit{Consigna del Trabajo
    Práctico N.º~3: Procesador RISC-V Pipeline con Debug por UART}, FCEFyN,
    UNC, 2026. \doc{consigna.md}{Texto de la consigna}.
  \item Cátedra de Arquitectura de Computadoras, \textit{Requisitos del TP3},
    FCEFyN, UNC, 2026. \doc{requirements.md}{Especificación de requisitos}.
  \item Proyecto TP3, \textit{Diseño FPGA consolidado}, documentación técnica.
    \doc{architecture/system_overview.md}{Diseño FPGA consolidado}.
  \item Proyecto TP3, \textit{Inventario de módulos y agrupación RTL},
    documentación técnica.
    \doc{architecture/module_inventory.md}{Inventario de módulos}.
  \item Proyecto TP3, \textit{Pipeline de la CPU}, documentación técnica.
    \doc{architecture/cpu/cpu_pipeline.md}{Pipeline de la CPU}.
  \item Proyecto TP3, \textit{Control del pipeline y hazards}, documentación
    técnica.
    \doc{architecture/cpu/control_and_hazards.md}{Control y hazards}.
  \item Proyecto TP3, \textit{Arquitectura de memoria}, documentación técnica.
    \doc{architecture/memory/memory_architecture.md}{Arquitectura de memoria}.
  \item Proyecto TP3, \textit{Control de ejecución y sesión}, documentación
    técnica.
    \doc{architecture/session/execution_control.md}{Control de ejecución y
    sesión}.
  \item Proyecto TP3, \textit{Protocolo UART}, documentación técnica.
    \doc{protocol.md}{Protocolo UART}.
  \item Proyecto TP3, \textit{Plan de verificación de hardware} y
    \doc{verification/coverage.md}{matriz de cobertura}, documentación técnica.
  \item Proyecto TP3, \textit{Manifiesto de casos de verificación}.
    \doc{verification/cases.json}{Manifiesto de casos}.
  \item Proyecto TP3, \textit{Entorno de compilación y simulación}.
    \doc{development_environment.md}{Entorno de desarrollo}.
  \item Proyecto TP3, \textit{Documentación de viewrv}.
    \doc{software/viewrv/README.md}{Índice de la aplicación}.
  \item Proyecto TP3, \textit{Aceptación Basys 3 de la issue N.º~17}.
    \acceptancedoc{Guía y resultados de la PR N.º~37},
    revisión \code{cb3270b19b27}, 2026.
  \item Proyecto TP3, \textit{Evidencia de la sesión física}.
    \evidencefile{resultados-por-id.csv}{Resumen por ID} y
    \evidencefile{tp3-issue-17-evidencia.zip}{capturas, bitstream y reportes},
    publicación \code{tp3-17-evidencia}, 2026.
  \item A. Waterman y K. Asanović (eds.), \textit{The RISC-V Instruction Set
    Manual, Volume I: User-Level ISA}, versión 2.2, RISC-V Foundation, 2017.
    \href{https://github.com/riscv/riscv-isa-manual/releases/tag/riscv-user-2.2}{Publicación oficial}.
    Versión adoptada en \doc{decisions/DEC-SYS-010.md}{DEC-SYS-010}.
  \item IEEE, \textit{Standard for SystemVerilog---Unified Hardware Design,
    Specification, and Verification Language}, IEEE 1800-2012.
    \href{https://standards.ieee.org/ieee/1800/4934/}{Ficha del estándar};
    versión de lenguaje seleccionada en la compilación RTL.
  \item Pele, \textit{Mermaid diagrams as themeable SVG}, obsidianmd/pele,
    licencia MIT. Disponible en \url{https://pele.run}.
  \item Digilent, \textit{Basys 3 Reference Manual}. Disponible en
    \url{https://digilent.com/reference/_media/reference/programmable-logic/basys-3/basys3_rm.pdf}.
  \item Digilent, \textit{Basys 3 Master XDC}. Disponible en
    \url{https://github.com/Digilent/digilent-xdc}.
  \item AMD, \textit{Vivado Design Suite User Guide: Using Constraints}, UG903.
    Disponible en \url{https://docs.amd.com/r/en-US/ug903-vivado-using-constraints}.
  \item AMD, \textit{Clocking Wizard v6.0 Product Guide}, PG065. Disponible en
    \url{https://docs.amd.com/r/en-US/pg065-clk-wiz}.
  \item AMD, \textit{Vivado Design Suite User Guide: Programming and Debugging},
    UG908. Disponible en
    \url{https://docs.amd.com/r/en-US/ug908-vivado-programming-debugging}.
  \item Icarus Verilog, \textit{Icarus Verilog Documentation}, versión 13.0.
  \item Verilator, \textit{Verilator User's Guide}, versión 5.020.
\end{itemize}
\endgroup
